% from ega4-8-fr.tex
\oldpage[IV-3]{25}
\[
\begin{tikzcd}
\varinjlim\mathfrak Q(\mathcal F_\lambda)
  \ar[r,"u_{\mathfrak Q}"] \ar[d]
  & \mathfrak Q(\mathcal F) \ar[d] \\
\varinjlim\mathfrak Q(\mathcal F'_\lambda)
  \ar[r,"u_{\mathfrak Q}"']
  & \mathfrak Q(\mathcal F')
\end{tikzcd}
\tag{8.5.10.2}\label{IV.8.5.10.2-fr}
\]
\end{env}

\begin{proposition}[8.5.11]
\label{IV.8.5.11-fr}
Supposons $S_0$ quasi-compact (resp. quasi-compact et quasi-séparé). Supposons
$\mathcal F_\lambda$ quasi-cohérent de type fini (resp. de présentation finie)
pour tout $\lambda$~; alors l'application canonique (8.5.10.1) est injective
(resp. bijective).
\end{proposition}

\begin{proof}
La première assertion résulte du lemme plus précis (8.5.9). Pour démontrer la
seconde, considérons un $\mathcal O_S$-Module quotient $\mathcal G$ de
$\mathcal F$ qui soit de présentation finie. Il résulte de (8.5.2, (ii)) qu'il
existe un $\lambda'\in L$ et un $\mathcal O_{S_{\lambda'}}$-Module
$\mathcal G_{\lambda'}$ quasi-cohérent de présentation finie tels que
$\mathcal G=u_{\lambda'}^*(\mathcal G_{\lambda'})$~; comme $L$ est filtrant, on
peut supposer $\lambda'=\lambda$ (en remplaçant $\lambda$ et $\lambda'$ par un
indice qui les majore). Considérons alors l'homomorphisme canonique
$p:\mathcal F\to\mathcal G$~; il résulte de (8.5.2, (i)) qu'il existe un
$\mu\geq\lambda$ et un homomorphisme
$p_\mu:\mathcal F_\mu\to\mathcal G_\mu$ tels que $p=u_\mu^*(p_\mu)$. En outre,
en vertu de (8.5.7), on peut supposer $\mu$ pris assez grand pour que $p_\mu$
soit surjectif, ce qui termine la démonstration.
\end{proof}

\subsection{Sous-préschémas de présentation finie d'une limite projective de préschémas}
\label{subsection:IV.8.6-fr}

\begin{env}[8.6.1]
\label{IV.8.6.1-fr}
Étant donné un préschéma $Y$, désignons dans ce numéro par
$\mathfrak{Spr}(Y)$ l'ensemble ordonné (I, 4.1.10) des sous-préschémas de $Y$
qui sont de présentation finie sur $Y$ (1.6.1), par $\mathfrak{Spro}(Y)$
(resp. $\mathfrak{Sprf}(Y)$) la partie de $\mathfrak{Spr}(Y)$ formée des
sous-préschémas induits sur des ouverts (resp. des sous-préschémas fermés) de
$Y$, de présentation finie sur $Y$~; il revient au même de dire qu'un
sous-préschéma $Z$ de $Y$ appartient à $\mathfrak{Spro}(Y)$ (resp.
$\mathfrak{Sprf}(Y)$) ou qu'il est induit sur un ouvert et que l'espace
sous-jacent $Z$ est rétrocompact dans $Y$ (resp. qu'il est fermé et que l'idéal
$\mathcal J$ de $\mathcal O_Y$ qui le définit est de type fini, ce qui signifie
aussi que $j_*(\mathcal O_Z)\in\mathfrak Q(\mathcal O_Y)$ si $j:Z\to Y$ est
l'injection canonique) (1.6.1 et 1.4.5).
\end{env}

\begin{env}[8.6.2]
\label{IV.8.6.2-fr}
Nous conservons toujours les notations de (8.2.2) et supposons que $S_0$ est
un des $S_\lambda$. Soit $Y_\lambda$ un sous-préschéma de $S_\lambda$~; alors
$Y_\mu=u_{\lambda\mu}^{-1}(Y_\lambda)$ (resp.
$Y=u_\lambda^{-1}(Y_\lambda)$) est un sous-préschéma de $S_\mu$ pour
$\mu\geq\lambda$ (resp. de $S$)~; il est induit sur un ouvert (resp. il est
fermé) si $Y_\lambda$ l'est (I, 4.3.2) et de présentation finie sur $S_\mu$
(resp. $S$) si $Y_\lambda$ est de présentation finie sur $S_\lambda$
(1.6.2, (iii)). Par suite
$(\mathfrak{Spr}(S_\lambda),u_{\lambda\mu}^{-1})$ (resp.
$(\mathfrak{Spro}(S_\lambda),u_{\lambda\mu}^{-1})$,
$(\mathfrak{Sprf}(S_\lambda),u_{\lambda\mu}^{-1})$) est un système inductif et
les applications
$u_\lambda^{-1}:\mathfrak{Spr}(S_\lambda)\to\mathfrak{Spr}(S)$ (resp.
$\mathfrak{Spro}(S_\lambda)\to\mathfrak{Spro}(S)$,
$\mathfrak{Sprf}(S_\lambda)\to\mathfrak{Sprf}(S)$) forment un système inductif
d'applications~; d'où, en passant à la limite inductive, des applications
canoniques
\begin{align}
  \varinjlim\mathfrak{Spr}(S_\lambda)
    &\longrightarrow\mathfrak{Spr}(S)
    \tag{8.6.2.1}\label{IV.8.6.2.1-fr}\\
  \varinjlim\mathfrak{Spro}(S_\lambda)
    &\longrightarrow\mathfrak{Spro}(S)
    \tag{8.6.2.2}\label{IV.8.6.2.2-fr}\\
  \varinjlim\mathfrak{Sprf}(S_\lambda)
    &\longrightarrow\mathfrak{Sprf}(S).
    \tag{8.6.2.3}\label{IV.8.6.2.3-fr}
\end{align}
\end{env}

